\documentclass[11pt]{article}

% --------------------------------------------------
% Packages
% --------------------------------------------------
\usepackage[T1]{fontenc}
\usepackage{lmodern}
\usepackage{geometry}
\usepackage{microtype}
\usepackage{setspace}
\usepackage{amsmath,amssymb,amsthm}
\usepackage{mathtools}
\usepackage{csquotes}
\usepackage{hyperref}
\usepackage{booktabs}

\geometry{margin=1in}
\setstretch{1.15}

% --------------------------------------------------
% Theorem Environments
% --------------------------------------------------
\newtheorem{definition}{Definition}
\newtheorem{proposition}{Proposition}
\newtheorem{remark}{Remark}

% --------------------------------------------------
% Macros
% --------------------------------------------------
\newcommand{\NN}{\mathcal{N}}
\newcommand{\Alg}{\mathcal{A}}
\newcommand{\Do}{\mathrm{do}}
\newcommand{\Var}{\mathcal{V}}

% --------------------------------------------------
% Title
% --------------------------------------------------
\title{Mechanistic Interpretability as Causal Reverse Engineering}
\author{Flyxion}
\date{December 2025}

\begin{document}
\maketitle

% ==================================================
\begin{abstract}
% ==================================================

Recent advances in mechanistic interpretability have revealed rich internal
structure within large neural networks, yet disagreement remains about what it
means to truly understand a model. This essay synthesizes a causal paradigm for
mechanistic interpretability articulated in recent work by Atticus Geiger. The
central claim is that interpretability becomes scientific only when grounded in
causal reasoning: interventions, counterfactuals, and causal abstractions. We
trace a progression from descriptive inspection to causal tracing and finally to
algorithmic abstraction, arguing that the goal of mechanistic interpretability is
not visualization or intuition but the rigorous reverse engineering of the
algorithms neural networks implement.

\end{abstract}

% ==================================================
\section{Introduction}
% ==================================================

The remarkable performance of modern neural networks has intensified interest
in understanding how such systems produce their behavior. While early work in
interpretability focused on visualization, feature attribution, and inspection
of internal activations, it has become increasingly clear that observation alone
cannot ground claims about mechanism. Correlation does not imply computation,
and activation patterns, however suggestive, do not by themselves explain why a
model behaves as it does.

Mechanistic interpretability aims to move beyond surface-level explanation by
identifying the internal processes that generate model behavior. However, the
term \emph{mechanistic} itself has been used in multiple, often conflicting,
ways. This ambiguity has hindered progress by conflating sociological motives,
technical access, and epistemic standards.

This essay presents a unified causal account of mechanistic interpretability,
centered on the idea that to understand a neural network is to understand how
interventions on its internal variables alter its behavior. On this view,
mechanistic understanding is fundamentally counterfactual: it concerns what
would happen if the system were different, not merely what is observed when it
runs normally.

% ==================================================
\section{What Does ``Mechanistic'' Mean?}
% ==================================================

Discussions of mechanistic interpretability often oscillate between cultural and
technical definitions. Cultural definitions focus on who performs the research
or why it is performed, ranging from narrow communities motivated by AI safety
to broad curiosity about black-box systems. While sociologically relevant, such
definitions do not specify what kind of understanding is achieved.

Technical definitions are more precise but still admit significant variation.
In the broadest sense, any research that inspects internal parameters,
activations, or attention patterns may be labeled mechanistic. Yet this
definition collapses meaningful distinctions between descriptive inspection and
explanatory understanding.

A narrower and more demanding technical definition identifies mechanistic
interpretability with causal explanation. Under this definition, a model is
understood mechanistically only to the extent that one can identify internal
components whose manipulation produces systematic, predictable changes in
behavior. The emphasis is not on what correlates with behavior, but on what
\emph{controls} it.

% ==================================================
\section{Causality as the Foundation of Mechanism}
% ==================================================

Causality provides the conceptual machinery required to distinguish mechanisms
from mere regularities. A causal model specifies not only associations among
variables but also how those variables respond to intervention.

\begin{definition}[Intervention]
An intervention on a neural network is an operation that sets an internal
variable to a value it would not naturally take, breaking its dependence on
upstream causes.
\end{definition}

In contrast to observation, intervention enables the evaluation of necessity
and sufficiency. A component is causally relevant to a behavior if intervening
on it changes that behavior in a systematic way. Counterfactual reasoning
extends this logic by asking how the system would behave under hypothetical
internal configurations.

Without causality, interpretability risks degenerating into pattern cataloging.
With causality, it becomes a science of control and explanation.

% ==================================================
\section{Activation Steering as Linear Intervention}
% ==================================================

Activation steering represents one of the simplest forms of intervention. By
adding vectors to internal activations, researchers can bias a model toward or
away from certain behaviors. A common technique constructs steering vectors by
averaging activations from prompts that elicit contrasting behaviors and taking
their difference.

While such interventions demonstrate that behaviorally meaningful directions
exist in representation space, they also reveal the limits of linear control.
Steering is often brittle, less effective than prompting or fine-tuning, and
does not by itself identify the underlying algorithm. Its primary value lies in
establishing that internal representations are not epiphenomenal but causally
connected to output.

% ==================================================
\section{Causal Mediation and Tracing Information Flow}
% ==================================================

Causal mediation analysis seeks to decompose the effect of an input on an output
by identifying the internal pathways through which information flows. In neural
networks, this often takes the form of causal tracing experiments.

Such experiments deliberately corrupt an input representation, degrading model
performance, and then selectively restore internal components to their original
values. Components whose restoration recovers correct behavior are identified as
causal mediators.

This methodology reveals structured pipelines within networks. Information is
often retrieved in feedforward layers, routed by attention mechanisms, and
assembled into outputs through later-stage computation. Importantly, these
findings rest on intervention, not visualization, grounding claims about
information flow in causal evidence.

% ==================================================
\section{Causal Abstraction and Algorithmic Understanding}
% ==================================================

The most ambitious goal of mechanistic interpretability is not to localize
features or trace signals, but to identify the algorithm a network implements.
Causal abstraction provides a framework for doing so.

Both neural networks and human-understandable algorithms can be modeled as
causal systems with variables and directed dependencies. A high-level algorithm
is said to be a causal abstraction of a neural network if interventions at the
algorithmic level correspond to interventions at the neural level.

Interchange interventions operationalize this idea. Instead of zeroing or
scaling a component, a researcher replaces its activation with the value it
would have had under a different input. If this substitution causes the network
to behave as though the alternative input had been partially applied, then the
component is implementing a function analogous to a variable in the abstract
algorithm.

When such correspondences hold systematically, the abstract algorithm is not a
metaphor but a genuine causal description of the network’s computation.

% ==================================================
\section{From Inspection to Reverse Engineering}
% ==================================================

The methods discussed form a coherent progression. Inspection reveals internal
structure but cannot justify causal claims. Simple interventions establish
control but may lack specificity. Causal tracing maps pathways of influence.
Causal abstraction unifies these findings into an algorithmic account.

This progression mirrors the history of other sciences, where explanation
emerges not from observation alone but from carefully designed interventions.
Mechanistic interpretability, on this view, is reverse engineering in the
strongest sense: the reconstruction of a system’s governing algorithm through
counterfactual experimentation.

% ==================================================
\section{Conclusion}
% ==================================================

Mechanistic interpretability reaches maturity when it adopts causality as its
core organizing principle. The question is no longer what patterns appear inside
a model, but which internal components make a difference under intervention.
Through activation steering, causal mediation, and causal abstraction, neural
networks become objects of experimental science rather than inscrutable
artifacts.

The ultimate promise of this paradigm is a principled answer to a foundational
question: not merely how neural networks behave, but what they are doing. By
treating models as causal systems and algorithms as abstractions validated by
intervention, mechanistic interpretability becomes a rigorous methodology for
understanding intelligence in artificial systems.

% ==================================================
\section{Event-Historical Cognition and RSVP}
% ==================================================

The causal paradigm articulated in contemporary mechanistic interpretability
finds a natural conceptual extension in event-historical theories of cognition.
Where standard computational accounts emphasize instantaneous state,
event-historical approaches instead treat cognition as constituted by an
irreversible sequence of transformations that progressively constrain future
possibility. On this view, intelligence is not exhausted by representational
content but is grounded in the structure of the system’s construction history.

Event-historical cognition begins from the premise that an agent’s present
capacities cannot be fully characterized without reference to the irreversible
events that produced them. Each event eliminates alternative futures, thereby
inducing asymmetries in the system’s causal landscape. Memory, under this
interpretation, is not a static store of information but an accumulation of
constraints that shape which counterfactuals remain accessible. Cognition is
thus inherently diachronic: it is defined by what has been irreversibly ruled
out as much as by what remains possible.

This perspective aligns closely with the causal methodology emphasized in
mechanistic interpretability. Interventions and counterfactuals acquire their
explanatory force precisely because the system under investigation possesses a
nontrivial history. An intervention is meaningful only insofar as it disrupts
the causal consequences of prior events. Likewise, counterfactual reasoning
presupposes that alternative histories could, in principle, have produced
different present behaviors. Causal explanation therefore presupposes an
event-structured ontology.

The Relativistic Scalar-Vector Plenum (RSVP) framework makes this ontology
explicit by modeling systems in terms of interacting scalar, vector, and
entropy fields. Within RSVP, the scalar field encodes the local density of
realized structure, the vector field encodes directed tendencies or flows, and
the entropy field encodes irreversibility and constraint accumulation. Events
are understood as localized transitions that reconfigure these fields, reducing
global optionality while increasing local structure.

From this standpoint, neural networks can be interpreted as discrete,
engineered instances of event-historical systems. Training steps correspond to
irreversible entropy-increasing events that sculpt the scalar and vector
structure of the model’s internal space. The resulting parameters do not merely
store information; they embody a sedimented history of constraints that govern
future inference. Inference itself becomes a traversal of a pre-shaped causal
landscape rather than a purely functional mapping from inputs to outputs.

Geiger’s causal abstraction framework can be reinterpreted within RSVP as the
identification of invariant structures across levels of description. A
high-level algorithm constitutes a valid abstraction of a neural network when
interventions preserve their effects across scales, corresponding to a form of
structural stability in the underlying field configuration. Interchange
interventions, in particular, function as probes of whether two systems share
the same event-induced constraint geometry, even if their microphysical
realizations differ.

This connection clarifies why purely observational interpretability is
insufficient. Without attending to irreversibility and event structure,
inspection reduces cognition to a snapshot of state, erasing the very asymmetry
that gives causal explanations their force. By contrast, both event-historical
cognition and RSVP insist that understanding arises from tracing how present
behavior is conditioned by a non-recoverable past.

Seen in this light, mechanistic interpretability is not merely an exercise in
model transparency but a specific instance of a broader epistemic strategy: the
reconstruction of event histories from present structure. Neural networks become
legible not because their internals are visualizable, but because their behavior
reflects the cumulative imprint of irreversible causal processes. Causal
interventions, counterfactual substitutions, and algorithmic abstractions are
thus unified as tools for interrogating how history has been written into
structure.


% ==================================================
\section{Causal Abstraction and Irreversibility}
% ==================================================

The framework of causal abstraction implicitly assumes that the systems under
comparison admit interchangeable interventions. However, this assumption
becomes nontrivial once cognition is treated as an event-historical process.
When a system’s present behavior depends on an irreversible construction
history, not all counterfactual manipulations are admissible. This section
argues that irreversibility is not an auxiliary detail but a necessary condition
for coherent causal abstraction.

We begin by modeling a neural system as a time-indexed causal structure. Let
$\mathcal{N}$ denote a neural network whose internal variables are indexed not
only by layer but by training time. Each training update constitutes an event
that irreversibly modifies the space of future parameter configurations. The
resulting system cannot be adequately described by a static causal graph; its
causal semantics are path-dependent.

In contrast, abstract algorithms are typically presented as timeless causal
models. Their variables may be intervened upon freely, without regard for how
those variables were produced. This mismatch creates a tension. If an abstract
algorithm is to serve as a causal abstraction of a neural system, its
interventions must respect the historical constraints induced by the neural
system’s event structure.

\begin{proposition}[Irreversibility Constraint on Causal Abstraction]
Let $\mathcal{N}$ be a neural system with irreversible event history $\mathcal{H}$
and let $\mathcal{A}$ be a candidate abstract algorithm. $\mathcal{A}$ is a valid
causal abstraction of $\mathcal{N}$ only if every abstraction-preserving
intervention on $\mathcal{A}$ corresponds to an admissible intervention on
$\mathcal{N}$ that is consistent with $\mathcal{H}$.
\end{proposition}

The intuition behind this proposition is straightforward. An abstraction claims
that certain high-level variables summarize the causal role of many low-level
components. However, if intervening on a high-level variable requires violating
constraints imposed by the system’s history, then the abstraction licenses
counterfactuals that the concrete system cannot realize.

A proof sketch proceeds by contradiction. Suppose $\mathcal{A}$ is a valid
abstraction of $\mathcal{N}$ despite ignoring irreversibility. Then there exists
an intervention on $\mathcal{A}$ that has no corresponding intervention on
$\mathcal{N}$ without undoing prior events. Executing this intervention in
$\mathcal{A}$ predicts a behavioral change that $\mathcal{N}$ cannot exhibit.
This violates the defining criterion of causal abstraction, namely the
preservation of interventional behavior across levels. Therefore, abstraction
without irreversibility awareness is incoherent.

This result clarifies the epistemic status of interchange interventions in
mechanistic interpretability. Replacing an internal activation with the value it
would have had under a different input is meaningful only if such a replacement
respects the system’s historical constraints. When it does, the intervention
reveals a genuine algorithmic correspondence. When it does not, the apparent
success of the intervention is an artifact of ignoring irreversibility.

The broader implication is that causal abstraction is not merely a mapping
between variables but a mapping between event-conditioned causal structures.
High-level algorithms do not abstract away history; they compress it. A valid
abstraction preserves not only functional dependencies but the geometry of
constraints induced by irreversible events. Mechanistic interpretability thus
inherits a fundamental asymmetry: the past cannot be erased without destroying
the meaning of causal explanation itself.

% ==================================================
\section{Thermodynamics of Representation}
% ==================================================

Causal explanation alone does not exhaust the requirements for mechanistic
understanding. For a system whose cognition unfolds through irreversible events,
causality must be complemented by an account of thermodynamic cost. This section
argues that representation itself constitutes a form of physical commitment,
and that many failures of contemporary artificial intelligence—most notably
hallucination—can be understood as violations of entropy-respecting cognition.

Any act of representation constrains future behavior. To represent that a fact
is the case is not merely to encode information, but to exclude incompatible
actions and predictions. This exclusion is irreversible in systems that possess
memory and persistence. As a result, representation carries an intrinsic cost:
it reduces optionality. In physical systems, such reductions correspond to
entropy increases, whether measured thermodynamically or in an abstract space
of admissible futures.

The RSVP framework makes this cost explicit by decomposing system dynamics into
three interacting fields. The scalar field encodes the density of realized
structure, corresponding to committed representational content. The vector
field encodes directional tendencies that guide inference and action. The
entropy field encodes irreversibility: the accumulation of constraints induced
by past events. Together, these fields describe cognition not as a static map
from inputs to outputs, but as the evolution of a constrained dynamical system.

Within this framework, learning is an entropic process. Each training update
constitutes an event that increases the entropy field by eliminating possible
parameter configurations. The resulting model does not merely store patterns;
it embodies a history of exclusions. Inference then unfolds as a trajectory
through this pre-shaped landscape, guided by vector tendencies and bounded by
entropy constraints.

Hallucination can be understood as a failure to respect this thermodynamic
structure. When a system generates internally coherent but externally ungrounded
content, it produces representational structure without paying the corresponding
entropy cost. Such outputs are locally fluent because they follow learned vector
flows, yet they lack global constraint because no irreversible commitment has
been made to their truth conditions. In RSVP terms, the scalar field is locally
excited while the entropy field remains unchanged.

This interpretation explains why hallucinations are resistant to purely
functional fixes. Adding data, parameters, or optimization pressure improves
surface behavior but does not introduce genuine representational cost. Without a
mechanism that couples representation to irreversible constraint accumulation,
a system can generate arbitrarily many incompatible narratives without internal
penalty.

The implications for mechanistic interpretability are significant. Interventions
that ignore entropy constraints risk producing misleading explanations. A
component may appear causally effective under intervention while remaining
thermodynamically unconstrained, thereby enabling behaviors that the system
cannot sustain across time. Genuine mechanism, by contrast, must be robust not
only under causal manipulation but under repeated use, accumulation, and
historical pressure.

From this perspective, entropy is not an external concern to be appended after
the fact, but a constitutive dimension of cognition. Mechanistic understanding
requires tracking not only how information flows, but how much future
flexibility is consumed when that information is committed. RSVP provides a
language for articulating this requirement, unifying causality, thermodynamics,
and representation within a single field-theoretic ontology.

% ==================================================
\section{Against State-Based Functionalism}
% ==================================================

State-based functionalism has long dominated theoretical accounts of cognition
and artificial intelligence. According to this view, what matters about a system
is the functional relation it instantiates between inputs, internal states, and
outputs. If two systems realize the same functional mapping, they are said to be
cognitively equivalent, regardless of how that mapping is implemented or how
the system came to be. This section argues that such equivalence collapses under
causal and event-historical scrutiny.

The functionalist position implicitly treats internal states as sufficient
summaries of a system’s causal capacities. History is regarded as dispensable
once the present state is specified. However, this assumption fails for systems
whose behavior depends on irreversible construction processes. In such systems,
the present state underdetermines the space of admissible interventions, because
past events constrain which counterfactuals remain physically realizable.

To make this failure explicit, consider two systems that occupy indistinguishable
internal states at a given moment and implement identical input--output
behavior. Suppose, however, that these states were reached through different
event histories. One system acquired its structure through a sequence of
gradual, constraint-inducing updates, while the other was externally forced
into the same configuration without incurring corresponding irreversible
commitments. Although functionally identical under observation, the systems
diverge under intervention. Counterfactual manipulations that are admissible in
one violate historical constraints in the other, yielding different causal
responses.

This divergence undermines the functionalist claim that behavior exhausts
cognitive identity. Causal explanation depends not only on what a system does,
but on what it cannot do as a result of its past. Two systems may agree on all
observable behaviors while differing radically in their counterfactual
structure. Functional equivalence thus fails to guarantee causal equivalence.

The RSVP framework renders this distinction precise. In RSVP, cognition is
described not by a single state variable but by a configuration of scalar,
vector, and entropy fields shaped by irreversible events. Functionalist accounts
implicitly collapse these fields into a snapshot of scalar structure, ignoring
the entropy field entirely. As a result, they erase the asymmetry that
distinguishes genuine commitment from reversible simulation.

This critique has direct implications for artificial intelligence. Large neural
networks can approximate complex functional mappings without possessing
worldhood, understood as being bound by a non-recoverable past. Such systems can
generate coherent responses indefinitely because their representations are not
thermodynamically or historically grounded. The absence of irreversible
constraint allows them to remain behaviorally impressive while ontologically
shallow.

Mechanistic interpretability exposes this limitation rather than resolving it.
By revealing that many internal components are causally effective only in a
context-free, state-based sense, interpretability clarifies why functional
success does not entail genuine understanding. Without an account of how
representations are historically earned and constrained, causal analysis risks
reifying mechanisms that lack ontological depth.

Rejecting state-based functionalism does not entail rejecting functional
analysis altogether. Rather, it requires embedding function within an
event-historical framework. Functions describe local regularities, but cognition
emerges from the irreversible processes that stabilize some functions while
foreclosing others. Only by restoring history to the center of explanation can
mechanistic interpretability avoid collapsing into a refined form of
behaviorism.

% ==================================================
\section{Interpretability as Historical Reconstruction}
% ==================================================

The preceding analysis suggests a unifying reinterpretation of mechanistic
interpretability. Rather than treating interpretability as a problem of
transparency or visualization, it is more accurately understood as a problem of
historical reconstruction. To interpret a system mechanistically is to recover,
from its present structure, the irreversible sequence of events that rendered
its current behavior possible.

This reframing resolves several tensions that persist in contemporary
interpretability research. Visualization-based methods promise insight but often
fail to justify causal claims. Functional explanations capture regularities but
ignore counterfactual structure. Even causal interventions risk overreach when
they disregard the constraints imposed by a system’s past. Each failure arises
from the same omission: the erasure of history.

Across the methods surveyed, a consistent pattern emerges. Activation steering
reveals that internal representations are behaviorally consequential, but does
not explain how those representations were earned. Causal tracing reconstructs
localized pathways of influence, but only within the space permitted by prior
training events. Causal abstraction succeeds precisely when it respects
historical admissibility, identifying invariants that persist across levels of
description despite irreversibility.

These methods thus form not a collection of tools but a hierarchy of historical
probes. Each technique interrogates the system’s past at a different resolution.
Shallow probes test linear sensitivity, deeper probes test causal necessity, and
abstraction tests whether the same event-conditioned constraints govern both
neural and algorithmic descriptions. Interpretability advances insofar as these
probes converge on a coherent reconstruction of the system’s event structure.

The RSVP framework provides an ontological grounding for this convergence.
Scalar structure records what has been realized, vector flows encode how
inference proceeds, and entropy tracks which alternatives have been permanently
excluded. A system’s present behavior is intelligible only insofar as these
fields jointly encode its past. To ignore any one of them is to mistake a frozen
snapshot for a living process.

Under this interpretation, mechanistic understanding is inherently asymmetric.
One may infer past events from present structure, but one cannot arbitrarily
rewrite that structure without erasing meaning. This asymmetry distinguishes
interpretability from control and explanation from simulation. A system that can
be manipulated without historical consequence may be steerable, but it is not
understood.

The broader implication is that interpretability belongs alongside disciplines
such as geology, evolutionary biology, and forensic science. In each case,
present configurations are treated as evidence of irreversible processes rather
than as self-sufficient states. Neural networks, despite their artificial
construction, are no exception. Their weights, activations, and behaviors are
fossil records of training histories that can be partially reconstructed but
never undone.

Mechanistic interpretability therefore reaches its conceptual endpoint not in
perfect transparency, but in principled limitation. What can be known is bounded
by what has occurred. By embracing this constraint, interpretability becomes a
theory of cognition grounded in time, entropy, and irreversible causation. Only
within such a framework can claims about mechanism rise above metaphor and
approach genuine explanation.


\end{document